\section{罪数}
\begin{enumerate}
  \item \zd{为了保障人权不能重复评价、重复处罚，为了保护法益不能遗漏评价、遗漏处罚}。
  \item 做题，\zd{先判断犯罪行为有几个}：
        \begin{enumerate}
          \item 1 个。择一重。
                \begin{enumerate}
                  \item \zd{继续犯}。重婚、非法拘禁。即使既遂，后续参与者也能构成承继的共同犯罪。
                  \item \zd{法条竞合}。一个行为多个同性质罪名。\al{合同诈骗，诈骗}。
                  \item \zd{想象竞合}。一个行为多个不同性质罪名。
                  \item 区分。触犯的必然性不同。触发合同诈骗罪必然触发诈骗罪。
                  \item \zd{加重犯}。基本犯 + 加重情节（重心）。
                        \begin{enumerate}
                          \item 情节。考试考\zd{基本犯和加重犯形态不一致}。
                                \begin{itemize}
                                  \item 都既遂，那就是正常情况，正常处理。
                                  \item 加重既遂，基本未遂/中止 \yjt 定加重犯的既遂，适用升格的法定刑，再按基本犯未遂/中止的规定从宽。
                                  \item 基本既遂，加重未遂/中止 \yjt 哪个重选哪个。
                                \end{itemize}
                          \item 结果。
                        \end{enumerate}
                \end{enumerate}
          \item 多个。
                \begin{enumerate}
                  \item 结合犯。法律特殊规定，没有就走正常的数罪并罚。
                        \begin{enumerate}
                          \item \zd{拐卖 + 强奸；绑架 + 杀人；绑架 + 故意伤害致人重伤死亡}；
                          \item 前罪 + 妨害公务罪。组织运送偷越国边境罪，走私贩卖运输制造毒品罪。
                          \item 交通肇事罪 + 因逃逸致人死亡。
                        \end{enumerate}
                  \item 连续犯。性质相同的行为连续实施。多次的两个功能。
                        \begin{enumerate}
                          \item 犯罪的成立条件 \yjt 多次盗窃、敲诈。
                          \item 法定升格条件 \yjt 多次抢劫。
                        \end{enumerate}
                  \item 吸收犯。重罪吸收轻罪。\zd{不遗漏评价}就大胆的吸收，否则数罪并罚。
                  \item 不可罚的事后行为。广义的吸收犯，只不过单独讨论。\zd{正常销赃}没必要定销赃罪，不具期待可能性。
                  \item 牵连犯。\al{想诈骗，然后伪造证件}。
                        \begin{enumerate}
                          \item 手段行为 + 目的行为有联系。
                          \item 常见、常发、常伴随。\al{想杀人，跑到公安偷枪}。
                        \end{enumerate}
                \end{enumerate}
        \end{enumerate}
  \item 再看罪名。罪名关系：
        \begin{enumerate}
          \item A 与非 A 的对立排斥。
          \item A 与 A + B 的包容评价。高位阶可以包容低位阶。法条竞合。
                \begin{itemize}
                  \item 故意杀人和故意伤害、强奸和强制猥亵、绑架和非法拘禁、拐卖儿童和拐骗儿童。
                  \item 抢劫罪 \yjt \yi 敲诈勒索；\er 抢夺罪 \yjt 盗窃罪 \yjt 侵占罪。
                \end{itemize}
          \item A 与 B 的中立关系。不搭嘎。故意杀人罪，盗窃罪。可以想象竞合。
          \item \zd{A 与 B 没关系，不必然意味 A 和 B 是对立排斥，还有可能是中立}。
        \end{enumerate}
\end{enumerate}
